%% ============================================================
%%  9.  CONCLUSION
%% ============================================================
\section{Conclusion}
\label{sec:conclusion}

We introduced Adaptive Influence Balancing (AIB), a framework that closes
the observe–control loop in multi-agent LLM debates by converting live
process diagnostics into adaptive intervention signals.
Our central finding is that correct minority hypotheses are prematurely
eliminated at an 8.2\% rate under a standard debate protocol, and that
this rate can be reduced by 61\% through diagnostic-triggered minority
protection, with a corresponding +5.5~pp gain in accuracy on MMLU-Pro.
A learned predictor further replaces hand-tuned thresholds with a
calibrated estimate of minority correctness, demonstrating that reasoning
quality and support trajectory are the dominant predictors.

Critically, not all interventions improve outcomes.
Adversarial loser resistance without quality filtering increases PER, and
categorical minority protection without diagnostic gating increases
influence asymmetry.
These failure modes reinforce the importance of
\emph{selective}, \emph{quality-gated}, \emph{diagnostics-informed}
intervention design.

\paragraph{Limitations.}
This study uses a single dataset (MMLU-Pro) and a single model family
(Qwen2.5).
The learned predictor is trained on 14B debates and evaluated on 7B
debates, introducing a model-size confound for Exp4.
No human preference evaluation has been conducted yet, though the study
protocol is designed and ready (\S\ref{sec:experiments}).

\paragraph{Future work.}
Immediate extensions include (i) multi-dataset validation on GPQA
(graduate-level scientific reasoning) and subjective policy debates,
(ii) heterogeneous agent teams (different model sizes or families),
(iii) the pre-registered human evaluation study (N=100 participants,
Prolific), and (iv) reinforcement-learning policies that learn when to
intervene from debate outcomes rather than from a separately trained
classifier.
